\begin{textsample}{NEG Dim 3 – 2009\_c6 – Score -3.00 – 2009\_c6/enade\_2009\_ci\_cont\_55d90a2674.txt}  \label{ex:f3_neg_045}
Capítulo 3 Análise Técnica da Prova Este capítulo tem por objetivo apresentar o desempenho dos estudantes de CIÊNCIAS CONTÁBEIS no ENADE/2009.
Para isso, foram calculadas as estatísticas básicas da prova como um todo, bem como as estatísticas das questões relacionadas à Formação Geral e ao Componente Específico.
Nas Tabelas, são apresentadas as seguintes estatísticas : total da população, de presentes, média, erro-padrão da média, desvio-padrão, mediana, nota máxima, nota mínima e coeficiente de assimetria ( ver definição em 1.4.3 ).
Nas estatísticas das questões discursivas, foram detalhadas também as notas nulas : em branco, por protesto ou por insuficiência.
As estatísticas apresentadas neste capítulo contemplaram, separadamente, os ingressantes, os concluintes e o total de estudantes, e foram calculadas tendo -se em vista as seguintes agregações : ( a ) Região Geográfica e País e ( b ) a Categoria Administrativa e a Organização Acadêmica.
As notas \textbf{brutas} das questões discursivas em cada uma das partes, Formação Geral/Componente Específico, foram calculadas a partir da média simples das notas de cada uma das questões correspondentes.
A primeira questão discursiva de Formação Geral foi desconsiderada e, portanto, a nota \textbf{bruta} referente às questões discursivas de Formação Geral foi igual à nota da questão mantida, a segunda.
No Componente Específico, a nota \textbf{bruta} correspondeu à média das notas nas três questões.
As notas \textbf{brutas} das questões objetivas de cada uma das partes, Formação Geral/Componente Específico, foram calculadas a partir dos acertos dos estudantes considerando -se uma regra de três para a conversão para a escala de 0 a 100, descontando -se aquelas questões que foram desconsideradas pela comissão ou que não passaram pelo teste do índice de discriminação ponto-bisserial ( ver seção 1.4.6 para a descrição do índice ).
As questões com índices de discriminação inferior a 0,2 foram excluídas do cômputo da nota do estudante.
As notas \textbf{brutas} de Formação Geral e do Componente Específico para cada um dos estudantes foram calculadas como médias ponderadas das suas notas nas questões objetivas e discursivas, ponderação diferente para cada parte : C\_FG = C\_FG\textasciicircum{}OBJ × 0,6 + C\_FG\textasciicircum{}DIS × 0,4 C\_CE = C\_CE\textasciicircum{}OBJ × 0,85 + C\_CE\textasciicircum{}DIS × 0,15, onde C\_FG é a nota \textbf{bruta} de Formação Geral, C\_FG\_OBJ é a nota \textbf{bruta} nas questões objetivas de Formação Geral, C\_FG\_DIS é a nota \textbf{bruta} nas questões discursivas de Formação Geral, C\_CE é a nota \textbf{bruta} do Componente Específico, C\_CE\_OBJ é a nota \textbf{bruta} nas questões objetivas do Componente Específico, e C\_CE\_DIS é a nota \textbf{bruta} nas questões discursivas do Componente Específico.
A nota \textbf{bruta} da prova, C\_GER, por sua vez, é também uma média ponderada da nota de cada uma das partes, Formação Geral/Componente Específico : C\_GER = 0,25 × C\_FG + 0,75 × C\_CE. 3.1 Estatísticas Básicas da Prova 3.1.1 Estatísticas Básicas Gerais A Tabela 3.1 apresenta as estatísticas básicas da prova por grupo de estudantes.
O total de estudantes inscritos para a prova foi de 94190.
Destes, 79,4 \% compareceram, sendo que a presença foi maior entre concluintes ( 90,4 \% ) do que entre ingressantes ( 72,7 \% ).
A média geral da prova foi de 30,41, sendo que os ingressantes obtiveram média mais baixa, de 27,45, que os concluintes, cuja média foi 34,38.
O desvio-padrão geral foi de 12,6, sendo o do grupo de ingressantes menor ( 10,9 ), do que o do grupo dos concluintes ( 13,5 ), aqueles possuíram uma distribuição mais homogênea nas notas.
A nota máxima foi 88,9, obtida por pelo menos um concluinte ao passo que a maior nota obtida por um ingressante foi 88,2.
Os coeficientes de assimetria para os dois grupos considerados, o de concluintes e o de ingressantes, bem como para a população como um todo foram todos positivos, indicando uma maior concentração à esquerda da distribuição e um maior espalhamento à direita.
Uma comparação entre o desempenho de ingressantes e concluintes pode ser realizada a partir dos dados apresentados no Gráfico 3.1.
Cumpre destacar que as notas dos concluintes, de um modo geral, foram mais elevadas do que as dos ingressantes, cuja distribuição de notas está à esquerda dos primeiros.
A ideia de melhor desempenho é reforçada pelas notas medianas dos dois grupos : 33,2 entre os concluintes e 26,3 entre os ingressantes.
Nas faixas de notas até 30 pontos, o percentual de ingressantes é maior que o dos concluintes, situação esta que se inverte para as faixas de notas superiores a 30 pontos.
Ainda assim, aproximadamente 3,3 \% dos ingressantes e 13,4 \% dos concluintes situaram -se nas faixas de notas superiores a 50,0 pontos.
O Gráfico 3.2 apresenta informações referentes ao desempenho geral de ingressantes e concluintes, comparando os resultados considerando a Região do País, Categoria Administrativa e Organização Acadêmica.
Levando -se em conta as notas médias dos estudantes em cada Região, observa -se que, nos dois grupos, há uma variação expressiva entre as Regiões.
As notas médias mais elevadas entre os concluintes ocorreram nas Regiões Sul ( 37,69 ) e Sudeste ( 34,81 ), ambas acima da média nacional ( 34,37 ).
Em relação aos ingressantes, as médias mais altas foram obtidas, também, nas Regiões Sul ( 28,51 ) e Sudeste ( 27,59 ) nas demais Regiões a média foi inferior à média nacional ( 27,45 ).
Na Região Norte foi obtida a média de notas mais baixa entre os ingressantes ( 25,59 ) e os concluintes ( 30,44 ) na prova.
Cursos nesta área foram oferecidos em instituições federais, estaduais, municipais e privadas.
Levando -se em conta os agrupamentos dos estudantes em Categorias Administrativas, observou -se que a pontuação média mais elevada entre os concluintes e ingressantes foi encontrada em instituições de origem federal ( respectivamente 42,97 e 34,65 ) – ambas acima da média nacional.
Tendo como foco as notas médias dos estudantes agrupados por Organização Acadêmica, notamos que universidades apresentaram as maiores médias, 36,99 para concluintes e 29,28 para ingressantes, ao passo que as piores notas médias foram observadas nas Faculdades, 32,14 para concluintes e 26,06 para ingressantes. 3.1.2 Estatísticas Básicas de Formação Geral A Tabela 3.2 apresenta as estatísticas básicas de Formação Geral da prova, por grupo de estudantes.
O total de estudantes inscritos para a prova foi de 94190.
Destes, 79,4 \% compareceram, sendo que a presença foi maior entre concluintes ( 90,4 \% ) do que entre ingressantes ( 72,7 \% ).
A média geral da parte da prova de Formação Geral foi de 38,77, sendo que os ingressantes obtiveram média mais baixa, 37,96 que os concluintes, cuja média foi 39,86.
O desvio-padrão geral foi de 19,6, sendo o do grupo de ingressantes 19,5, e o do grupo dos concluintes 19,7, aqueles possuíram uma distribuição mais homogênea nas notas.
A nota máxima foi 100, obtida por estudantes de ambos os grupos.
Para esta parte, assim como para a nota geral, os coeficientes de assimetria são positivos para todos os grupos.
Como, porém, os valores são pequenos, estes indicaram uma ligeira inclinação para a esquerda.
O Gráfico 3.3 permite a comparação entre concluintes e ingressantes relativamente ao desempenho na parte da prova que aborda a Formação Geral.
Nesta parte, em se excluindo as notas nulas, tal como na prova como um todo, também os concluintes apresentaram desempenho superior, embora neste perceptível um maior equilíbrio.
A mediana dos concluintes foi igual à dos ingressantes, enquanto a média dos concluintes foi maior do que a dos ingressantes.
Dos ingressantes, 27,1 \% alcançaram nota superior a 50 ; dentre os concluintes, o percentual dos estudantes nesta faixa de notas foi muito semelhante e igual a 30,3 \%.
No Gráfico 3.4 são apresentadas as informações referentes ao desempenho dos ingressantes e dos concluintes na parte de Formação Geral, em diferentes agregações como Região, Categoria Administrativa e Organização Acadêmica.
Considerando -se as notas médias dos estudantes segundo cada Região, observou -se que, tanto no grupo dos ingressantes quanto no grupo dos concluintes, houve diferenças marcantes.
As médias variam, entre os concluintes, de 34,68 na Região Norte a 44,18 na Região Sul.
Entre os ingressantes, os valores foram de 33,27 na Região Norte a 40,56 na Região Sul.
Tendo como foco as Categorias Administrativas, as maiores notas médias foram das instituições federais, tanto para concluintes ( 49,33 ) como para ingressantes ( 47,88 ), enquanto as das instituições privadas foram menores tanto para concluintes ( 38,45 ) quanto para ingressantes ( 36,40 ).
Note que, com este recorte, ingressantes apresentaram médias de notas superiores aos dos concluintes nas instituições estaduais.
Os estudantes das instituições privadas se apresentaram mais homogêneos do que os das demais instituições.
Considerando -se o tipo de Organização Acadêmica, Universidades apresentaram os maiores valores para os concluintes e para os ingressantes.
Já os menores valores foram encontrados nas Faculdades, também para ambos os grupos de estudantes.
Os Centros Federais de Educação Tecnológica e os Institutos de Educação, Ciência e Tecnologia não estavam representados nesta área. 3.1.3 Estatísticas Básicas do Componente Específico A Tabela 3.3 apresenta as estatísticas básicas referentes ao Componente Específico da prova, tendo sua média geral, de 27,68, inferior ao desempenho da prova de CIÊNCIAS CONTÁBEIS como um todo, cuja média foi 30,41, como apresentado na Tabela 3.1, que por sua vez foi inferior ao desempenho na Formação Geral, que teve média igual a 38,77, conforme mostrado na Tabela 3.2.
Os concluintes apresentaram um desempenho médio ( 32,60 ) superior ao dos ingressantes ( 24,00 ), o oposto do ocorrido para a nota da Formação Geral, havendo ainda indícios de maior homogeneidade entre as notas médias dos ingressantes do que entre as dos concluintes, como sugeriu o desvio padrão, menor no grupo dos ingressantes do que no dos concluintes.
As notas máximas obtidas por concluintes e ingressantes foram 94,4 e 85,5, respectivamente.
As distribuições de notas para cada grupo se apresentaram com uma inclinação para notas menores, como evidenciado pelos coeficientes de assimetria positivos.
A situação de assimetria é confirmada pelo Gráfico 3.5 que proporciona uma comparação do desempenho entre concluintes e ingressantes em relação ao Componente Específico da prova.
Claramente a distribuição das notas médias dos ingressantes, à direita daquela dos concluintes, evidenciou o percentual mais alto de notas médias nas faixas mais baixas para este grupo, que como já comentado, incluiu uma maior proporção de zeros entre os ingressantes.
As medianas das distribuições das notas dos ingressantes e dos concluintes foram 22,6 e 30,6, respectivamente, enfatizando o melhor desempenho dos concluintes nesta parte da prova.
Assim como os Gráficos 3.2 e 3.4, o Gráfico 3.6 apresenta uma comparação dos resultados em relação às Regiões, Categoria Administrativa e Organização Acadêmica, agora levando em conta o desempenho dos ingressantes e dos concluintes no Componente Específico da prova.
Tendo como foco as notas médias dos estudantes, em determinada Região, observa -se que, entre os concluintes, as mais elevadas foram encontradas nas Regiões Sul ( 35,58 ) e Sudeste ( 32,99 ), ambas superiores à média nacional ( 32,59 ).
Entre os ingressantes, as maiores notas ocorreram também nas mesmas Regiões : Sul ( 24,54 ) e Sudeste ( 24,09 ), ambas acima da média nacional ( 23,99 ).
Novamente, para a análise segundo a Categoria Administrativa, todas as entidades se encontraram representadas, instituições federais, estaduais, municipais e privadas.
As médias mais elevadas, tanto entre os concluintes como entre os ingressantes, foram encontradas em instituições de origem federal ( 40,90 entre os concluintes e 30,29 entre os ingressantes ) enquanto nas instituições privadas a média para os concluintes foi de 31,57 e nas municipais foi de 22,76 para os ingressantes.
Quanto à Organização Acadêmica, só as universidades, centros universitários e faculdades foram representadas com este curso.
Considerando -se o tipo de Organização Acadêmica, as universidades apresentaram os maiores valores, seguidas pelos centros universitários e pelas faculdades.
Assim como os Gráficos 3.2 e 3.4, o Gráfico 3.6 apresenta uma comparação dos resultados em relação às Regiões, Categoria Administrativa e Organização Acadêmica, agora levando em conta o desempenho dos ingressantes e dos concluintes no Componente Específico da prova.
Tendo como foco as notas médias dos estudantes, em determinada Região, observa -se que, entre os concluintes, as mais elevadas foram encontradas nas Regiões Sul ( 35,58 ) e Sudeste ( 32,99 ), ambas superiores à média nacional ( 32,59 ).
Entre os ingressantes, as maiores notas ocorreram também nas mesmas Regiões : Sul ( 24,54 ) e Sudeste ( 24,09 ), ambas acima da média nacional ( 23,99 ).
Novamente, para a análise segundo a Categoria Administrativa, todas as entidades se encontraram representadas, instituições federais, estaduais, municipais e privadas.
As médias mais elevadas, tanto entre os concluintes como entre os ingressantes, foram encontradas em instituições de origem federal ( 40,90 entre os concluintes e 30,29 entre os ingressantes ) enquanto nas instituições privadas a média para os concluintes foi de 31,57 e nas municipais foi de 22,76 para os ingressantes.
Quanto à Organização Acadêmica, só as universidades, centros universitários e faculdades foram representadas com este curso.
Considerando -se o tipo de Organização Acadêmica, as universidades apresentaram os maiores valores, seguidas pelos centros universitários e pelas faculdades. 3.2 Análise das Questões Objetivas 3.2.1 Formação Geral A Tabela 3.4 apresenta as estatísticas básicas relativas às oito questões objetivas da Formação Geral dos estudantes.
Adiantamos que a primeira questão foi desconsiderada pelo critério da correlação ponto-bisserial.
Como pode ser observado na Tabela 3.4, os estudantes concluintes obtiveram desempenho superior, com nota média de 51,03, em relação aos estudantes ingressantes cuja nota média foi 49,22.
A variabilidade entre os concluintes foi menor do que entre os ingressantes, como indicaram os desvios-padrão de 21,8, para o grupo de concluintes e 22,1 para o grupo dos ingressantes.
A amplitude total das notas foi a mesma para ambos os grupos : nota mínima de 0 e máxima de 100.
A mediana para ingressantes e concluintes foi 42,8 e 57,1, respectivamente.
A assimetria, medida pelo coeficiente de assimetria ( skewness ) foi igual a zero para ambos os grupos, indicando uma simetria da distribuição das notas.
A Tabela 3.5 apresenta a classificação das questões objetivas da Formação Geral, segundo o índice de facilidade ( questões em vermelho foram desconsideradas pelo critério ponto-bisserial ).
Como as questões foram comuns a todas as áreas, a tabela apresenta a classificação considerando -se os concluintes de todas as áreas e aqueles do curso de CIÊNCIAS CONTÁBEIS.
De acordo com os índices obtidos, as questões objetivas da prova foram assim avaliadas para a área de CIÊNCIAS CONTÁBEIS : das oito questões, duas tiveram entre 60 \%, inclusive, e 85 \% de acertos e foram classificadas como fáceis ( somente uma para o total de concluintes ) ; outras duas foram consideradas medianas, alcançando entre 40 \%, inclusive, e 60 \% de acertos ( três para o total de concluintes ) ; três obtiveram entre 15 \%, inclusive, e 40 \% de acertos, sendo consideradas difíceis ( três para o total de concluintes ).
Por fim, uma foi classificada como muito difícil tento obtido percentual menor do que 15 \% de acertos ( uma para o total de concluintes ), coincidentemente a questão desconsiderada pelo critério ponto-bisserial.
Nenhuma questão foi classificada na categoria muito fácil.
Em linhas gerais podemos dizer que os concluintes de CIÊNCIAS CONTÁBEIS tiveram um melhor desempenho em Formação Geral em relação ao total da população de concluintes de todos os cursos.
A Tabela 3.6 descreve os resultados das análises das questões objetivas relativas à Formação Geral, segundo o poder de discriminação, utilizando -se para tal o índice de discriminação ponto-bisserial.
Como as questões foram comuns a todas as áreas, a tabela apresenta a classificação considerando -se os concluintes de todas as áreas e aqueles do curso de CIÊNCIAS CONTÁBEIS.
A desconsideração da questão segundo este critério considerou o universo de todos os concluintes.
Nenhuma questão apresentou índice igual ou acima de 0,40 ( nenhuma para a população de concluintes como um todo ), uma foi classificada como boa com índice entre 0,30, inclusive, e 0,40 ( três para os concluintes como um todo ) e cinco foram classificadas como médias, com índice entre 0,20, inclusive, e 0,30 ( quatro para os concluintes como um todo ).
Duas das questões foram classificadas como fracas ( uma para a população de concluintes como um todo ) sendo desconsiderada pelo critério ponto-bisserial.
Na sequência, no Quadro 3.1, estão especificadas, por questão, as habilidades aferidas pelas questões de Formação Geral da prova de CIÊNCIAS CONTÁBEIS, bem como os resultados da análise relativa aos índices de facilidade e de discriminação de cada questão.
O índice de facilidade para a área de CIÊNCIAS CONTÁBEIS teve uma variação considerável, oscilando entre 0,10 e 0,80 e o de discriminação ficou entre 0,03 e 0,37.
A questão 7, que aferiu as habilidades de “ Analisar e criticar informações ” ; ” Extrair conclusões por indução e/ou dedução ”, ” Estabelecer relações, comparações e contrastes em diferentes situações ”, “ Detectar contradições ”, “ Fazer escolhas \textbf{valorativas} avaliando consequências ”, “ Questionar a realidade ” e “ Argumentar coerentemente ” foi a que apresentou, para a população de concluintes de CIÊNCIAS CONTÁBEIS, maior poder discriminatório com índice de 0,37.
A questão mais fácil, com 80 \% de acertos, foi a de número 8, que aferiu também, além das já listadas anteriormente, a habilidade de “ Ler e interpretar textos ”.
De outra parte, a questão com menor índice de facilidade e de discriminação em Formação Geral foi a questão 1, que foi desconsiderada, e que deveria aferir as habilidades de “ Ler e interpretar textos ”, “ Analisar e criticar informações ” ; ” Extrair conclusões por indução e/ou dedução ”, ” Estabelecer relações, comparações e contrastes em diferentes situações ”, “ Detectar contradições ”, “ Fazer escolhas \textbf{valorativas} avaliando consequências ”, “ Questionar a realidade ” e “ Argumentar coerentemente ”.
O Gráfico 3.7 mostra, a título de ilustração, o comportamento das respostas à questão 3 de Formação Geral para os estudantes concluintes das áreas de CIÊNCIAS CONTÁBEIS como função da nota geral na prova4.
Trata -se de uma questão considerada média para os concluintes de CIÊNCIAS CONTÁBEIS e média para o conjunto de concluintes de todas as áreas.
O poder de discriminação foi médio para os estudantes da área e médio para o conjunto de todos os concluintes.
Neste Gráfico, cada uma das cinco curvas representou o percentual de respostas em determinada alternativa da questão, em função do escore total dos estudantes concluintes da área de CIÊNCIAS CONTÁBEIS.
A curva roxa corresponde à alternativa E, a correta para este item.
Assim, observa -se que entre os estudantes com notas mais baixas – desempenho geral mais fraco na prova como um todo –, ocorreu, além da escolha da alternativa certa, a escolha de alternativas incorretas, especialmente a alternativa C – curva vermelha, seguida da alternativa D – curva em verde escuro.
A proporção de estudantes que selecionaram a resposta correta, aumentou com o desempenho deste grupo de estudantes chegando ao acerto de 100 \% dos estudantes com notas pelo menos 87,5.
Esta análise permitiu verificar como a questão discriminou entre os dois grupos de desempenho de acordo com o índice obtido.
Os Gráficos relativos às demais questões de Formação Geral constam do Anexo I. 3.2.2 Componente Específico A Tabela 3.7 apresenta as estatísticas básicas em relação às questões objetivas do Componente Específico.
Como pode ser observado, os estudantes concluintes obtiveram desempenho superior, com uma média de 36,96 pontos, em relação aos estudantes ingressantes com uma média de 27,29 pontos.
Os desvios-padrão indicaram que a variabilidade entre os concluintes foi superior ( 16,6 ) do que entre os ingressantes ( 13,0 ).
Os coeficientes de assimetria para os dois grupos foram positivos indicando uma concentração à esquerda e espalhamento à direita.
A assimetria foi, porém, ligeiramente maior entre os ingressantes.
A análise das questões objetivas do Componente Específico em CIÊNCIAS CONTÁBEIS, em relação ao grau de facilidade da prova ( ver Tabela 3.8 ), mostrou que, das 27 questões, 19 encontraram -se na categoria difícil ( entre 15 \% e 40 \%, não incluso, de acertos ), 3 foram de nível médio ( entre 40 \% e 60 \%, não incluso, de acertos ), nenhuma questão está na categoria fácil ( entre 60 \% e 85 \%, não incluso, de acertos ) e 1 teve a classificação muito fácil ( a partir de 85 \% de acertos ), 2 questões tiveram percentual de acertos menor do que 15 \% ( classificação muito difícil ) e 2 questões foram desconsideradas pela Comissão Assessora da Área de CIÊNCIAS CONTÁBEIS Das 27 questões objetivas relativas ao Componente Específico ( ver Tabela 3.9 ), 5 obtiveram o índice de discriminação muito bom, 4 obtiveram o índice de discriminação bom e 6 tiveram nível de discriminação médio.
O índice fraco que eliminaria a questão, foi alcançado por 10 questões e 2 foram desconsideradas.
Na sequência, no Quadro 3.2 estão especificadas, por questão, as habilidades aferidas pelas questões do Componente Específico da prova, bem como os resultados da análise relativa aos índices de facilidade e de discriminação de cada questão.
Das 27 questões objetivas relativas ao Componente Específico ( ver Tabela 3.9 ), 5 obtiveram o índice de discriminação muito bom, 4 obtiveram o índice de discriminação bom e 6 tiveram nível de discriminação médio.
O índice fraco que eliminaria a questão, foi alcançado por 10 questões e 2 foram desconsideradas.
Na sequência, no Quadro 3.2 estão especificadas, por questão, as habilidades aferidas pelas questões do Componente Específico da prova, bem como os resultados da análise relativa aos índices de facilidade e de discriminação de cada questão.
O Gráfico 3.8 mostra, a título de ilustração, o comportamento da questão 13 do Componente Específico para os estudantes concluintes da área de CIÊNCIAS CONTÁBEIS como função da nota geral na prova.
Trata -se de uma questão considerada difícil para os concluintes de CIÊNCIAS CONTÁBEIS.
O poder de discriminação foi muito bom para os estudantes da área.
Neste Gráfico, cada uma das cinco curvas, representa o percentual de respostas em determinada alternativa da questão, em função da nota total dos estudantes concluintes de CIÊNCIAS CONTÁBEIS.
A curva verde claro correspondeu à alternativa B, a correta para esta questão.
Assim, observou -se que entre os estudantes com escores mais baixos – desempenho geral mais fraco na prova como um todo –, houve a escolha de alternativas incorretas, como a alternativa E – curva roxa e a alternativa D – curva verde escura.
A proporção de estudantes que selecionaram a resposta correta aumentou com o desempenho deste grupo de estudantes chegando ao acerto de 100 \% dos estudantes com notas pelo menos 87,5.
Esta análise permitiu verificar como a questão discriminou entre os grupos de desempenho de acordo com o índice obtido.
Os Gráficos relativos às demais questões do Componente Específico constam do Anexo I. 3.3 Análise das Questões Discursivas 3.3.1 Formação Geral A Tabela 3.10 apresenta as estatísticas básicas da prova por grupo de estudantes.
A análise dos resultados de desempenho dos estudantes nas questões discursivas que são apresentados na tabela 3.10 e no gráfico 3.9, mostrou que as notas foram mais baixas no conjunto dessas questões que no das questões objetivas.
O total de estudantes inscritos para a prova foi de 94190.
Destes, 79,4 \% compareceram, sendo que a presença foi maior entre concluintes ( 90,4 \% ) do que entre ingressantes ( 72,7 \% ).
A média geral da prova foi de 22,04, sendo que os ingressantes obtiveram média mais baixa, de 21,16, que os concluintes, cuja média foi 23,22.
O desvio-padrão geral foi de 31,6, sendo o do grupo de ingressantes menor ( 31,0 ), em relação ao grupo dos concluintes ( 32,3 ), que possuíram uma distribuição mais homogênea nas notas.
A nota máxima foi 100,0, atingida por concluintes e ingressantes.
Os coeficientes de assimetria para os dois grupos considerados, o de concluintes e o de ingressantes, bem como para a população como um todo foram todos positivos, indicando uma maior concentração à esquerda da distribuição.
Para todos os grupos, mais da metade dos estudantes presentes obtiveram nota zero nas questões discursivas de Formação Geral.
Das 46694 provas com nota zero, 1.559 ( 3,3 \% ) continham um protesto, 14.218 ( 30,4 \% ) estavam em branco e 30917 ( 66,2 \% ) foram respondidas, porém de forma inadequada.
Daqueles que tiveram nota zero, a maioria foi devido à resposta inadequada – 66,2 \% para ingressantes e concluintes.
As notas zero devido a protesto representaram 3,2 \% entre concluintes e 3,4 \% entre ingressantes.
Já a proporção de notas zero decorrentes de provas em branco, foi 30,6 \% entre concluintes e 30,4 \% entre ingressantes.
Uma comparação entre o desempenho de ingressantes e concluintes pode ser realizada a partir dos dados apresentados no Gráfico 3.9.
As notas dos concluintes, de um modo geral, foram mais elevadas do que as dos ingressantes, cuja distribuição de notas está à direita dos primeiros.
A ideia de melhor desempenho foi reforçada pela proporção de notas nulas.
Na sequência, discorrer -se-á, especificamente, sobre uma das questões discursivas de Formação Geral, a número 10 já que a número 9 foi desconsiderada.
Em especial, serão analisados os conteúdos englobados na questão e os resultados obtidos pelos estudantes.
Ao final, são apresentados os comentários da banca corretora a respeito do observado nas respostas dos estudantes, suas impressões e conclusões.
Cumpre esclarecer que, tendo em vista que as questões discursivas de Formação Geral foram padronizadas, ou seja, constam de todas as provas, os comentários da banca foram os mesmos para todas as áreas acadêmicas, eis que direcionados a todos os estudantes que participaram do ENADE/2009. 3.3.1.1 Análise da Questão Discursiva 9 de Formação Geral DESCONSIDERADA 3.3.1.2 Análise da Questão Discursiva 10 de Formação Geral A Questão Discursiva nº 10 buscou avaliar as habilidades : • Ler e interpretar textos ; • Analisar e criticar informações ; • Extrair conclusões por indução e/ou dedução ; • Estabelecer relações, comparações e contrastes em diferentes situações ; • Fazer escolhas \textbf{valorativas} avaliando consequências ; • Questionar a realidade ; • Argumentar coerentemente.
O Quadro 3.3 apresenta as habilidades/competências, temas predominantes e estatísticas básicas da questão 10 de Formação Geral por grupo de estudantes.
Cumpre notar que como a questão 9 foi desconsiderada, a nota referente às questões discursivas da Formação Geral é idêntica à nota da questão 10.
A média geral desta questão foi de 22,04, sendo que os ingressantes obtiveram média mais baixa, de 21,16, que os concluintes, cuja média foi 23,22.
O desvio-padrão geral foi de 31,6, sendo o do grupo de ingressantes menor ( 31,0 ), do que o do grupo dos concluintes ( 32,3 ), aqueles que possuíram uma distribuição mais homogênea nas notas.
A nota máxima foi 100,0, atingida por concluintes e ingressantes.
Os coeficientes de assimetria para os dois grupos considerados, o de concluintes e o de ingressantes, bem como para a população como um todo foram positivos, indicando uma maior concentração à esquerda da distribuição.
Para todos os grupos, mais da metade dos estudantes presentes tiveram nota zero nesta questão discursiva de Formação Geral.
Das 46694 provas com nota zero, 1.559 ( 3,3 \% ) continham um protesto, 14218 ( 30,4 \% ) estavam em branco e 30917 ( 66,3 \% ) foram respondidas, porém de forma inadequada.
Daqueles que tiveram nota zero, a maioria foi devido à resposta inadequada – 66,2 \% para ingressantes e concluintes.
As notas zero devido a protesto representaram 3,2 \% entre concluintes e 3,4 \% entre ingressantes.
Já a proporção de notas zero decorrentes de provas em branco, foi 30,6 \% entre concluintes e 30,4 \% entre ingressantes.
Uma comparação entre o desempenho de ingressantes e concluintes para esta questão foi exatamente igual à já realizada para a nota discursiva de Formação Geral a partir dos dados apresentados no Gráfico 3.9. 3.3.1.3 Comentários sobre as respostas à Questão Discursiva 10 Apresentação A questão 10 de Formação Geral apresentou um questionamento capaz de analisar a habilidade do estudante em reconhecer o que é um direito social, optar por um direito social vigente que possa diminuir as desigualdades de renda familiar no Brasil e discutir os argumentos em favor do direito social escolhido como suporte à resposta.
Para análise das respostas apresentadas pelos estudantes, houve uma leitura inicial bastante criteriosa.
Foram considerados para efeito de pontuação, os seguintes argumentos : • Acesso à educação pública, gratuita e de qualidade, garantindo ao cidadão melhores condições de vida ; • Assistência à saúde, em seu contexto mais amplo, garantindo os cuidados essenciais de saúde, baseados em técnicas apropriadas ao alcance de todas as pessoas, contando com a participação da população ; • Permanência do estudante na escola, em todos os níveis escolares ( da educação infantil à educação superior, possibilitando a qualificação profissional ), o que garante melhores condições de trabalho e remuneração.
Considerou -se a qualificação e o desempenho docente, desenvolvendo nos estudantes melhores expectativas de aprendizagem e interesse pelo conteúdo da escola.
Os professores desenvolvendo metodologias adequadas e conhecimento pleno quanto a questões de aprendizagem ; • Condições dignas de trabalho, com remuneração que garanta qualidade de vida.
Foram consideradas as condições de trabalho, a preparação profissional para o trabalho e os cursos de aperfeiçoamento, oportunizando momentos de aprendizagem e desenvolvimento, a partir de situações desencadeadoras ; • Ser proprietário do imóvel em que reside, reduzindo os gastos com aluguel, garantindo melhores condições de vida, equilíbrio financeiro e construção do bem-estar familiar.
Desempenho dos estudantes na questão É importante salientar que havia respostas que desenvolveram adequadamente o questionamento apresentado, com fundamentação baseada em argumentos coerentes.
Contudo, em um número considerável de respostas, havia apenas um argumento e não dois, conforme fora solicitado.
Muitos estudantes informam sobre problemas da educação, de um modo geral.
Algumas observações mencionadas pelos corretores durante o processo de correção foram : • há preocupação com uma educação de qualidade, com sugestões de cursos profissionalizantes e métodos preparatórios de aulas. • a educação é revelada como de baixa qualidade, exaltando -se treinamentos e cursos de atualização para a classe docente. • a saúde como essência para a vida.
O SUS é colocado como parte integrante de direito social, garantindo às pessoas fácil acesso. • oportunidade de emprego como forma de ascensão social, direito, cidadania. • os políticos brasileiros como causadores do despreparo educacional.
As respostas analisadas apresentaram norma culta, coesão e coerência razoáveis, que garantiram a compreensão da mensagem desenvolvida pelos estudantes.
Os erros ortográficos foram poucos e não comprometeram a discussão sobre o assunto e sobre os pontos de vista abordados.
Quanto à coerência, houve um desvio em relação ao objetivo da questão, já que alguns estudantes, através dos seus argumentos retrataram reivindicações, reclamações e situações de problemas enfrentados pelos brasileiros, citando golpes políticos.
Houve erros quanto à concordância verbal e nominal.
Houve, também, dificuldade em utilizar os conectores corretamente.
Alguns estudantes apresentaram em suas respostas redundância e prolixidade.
O desempenho foi considerado satisfatório, pois a maioria dos estudantes soube responder ao questionamento.
Houve domínio do conteúdo, pois as ideias, os argumentos e a \textbf{originalidade} corresponderam ao esperado. 3.3.2 Componente Específico Nas questões discursivas do Componente Específico – questões de números 38, 39 e 40 –, a diferença de desempenho entre os grupos não foi muito acentuada, inclusive menor do que para as questões discursivas de Formação Geral.
A Tabela 3.11 contém as estatísticas da média destas questões.
Observa -se que no grupo dos concluintes a nota média alcançada foi 8,03, enquanto que no grupo dos ingressantes foi de 5,49, ambas menores do que as obtidas nas questões discursivas de Formação Geral, que foram de 21,16 para os ingressantes e 23,22 para os concluintes ( Tabela 3.10 ).
As médias foram também menores do que as obtidas para as questões objetivas desse componente ( Tabela 3.7 ) : respectivamente 27,29 e 36,96 para ingressantes e concluintes.
A análise de cada uma destas questões será feita nas subseções seguintes.
Os valores positivos para a medida de assimetria indicaram uma distribuição com concentração à esquerda.
Dos 74747 presentes, 64,9 \% obtiveram nota zero.
A tabela 3.12 apresenta um detalhamento da situação destes alunos.
Menos de 0,05 \% não fez nenhuma das 3 questões.
Este limite serve para os concluintes ( 0,0 \% ) e ingressantes ( 0,0 \% ).
A combinação de questões em branco e preenchidas correspondeu a 98,2 \% das provas com nota zero na parte discursiva do Componente Específico.
Nenhuma das provas teve as 3 questões com protesto.
O Gráfico 3.10 apresenta a distribuição das notas nas questões discursivas do Componente Específico por grupo de estudantes.
A moda se situou no primeiro intervalo, aquele com valores entre 0 e 10 incluindo os extremos do intervalo.
A impressão de assimetria informada pelos coeficientes de assimetria positivos ( Tabela 3.11 ) foi reforçada no gráfico pela concentração no primeiro intervalo. 3.3.2.1 Análise da Questão Discursiva 38 do Componente Específico DESCONSIDERADA 3.3.2.2 Análise da Questão Discursiva 39 do Componente Específico A questão 39, cujas habilidades aferidas encontram -se descritas no Quadro 3.4, que ainda contém outras informações, seguido da Tabela 3.13 e do Gráfico 3.11, foi aquela em que os grupos de estudantes observados obtiveram o melhor desempenho entre as questões discursivas do Componente Específico.
A nota média do grupo dos ingressantes foi de 9,81, e a dos concluintes 16,49, bem melhor entre os concluintes.
Isto ocorre, principalmente pelo grande percentual de notas nulas : 72,5 \% entre os ingressantes e 64 \% entre os concluintes.
Das provas com notas nulas nesta questão, uma grande proporção foi deixada em branco : 73,2 \% entre os ingressantes e 73,2 \% entre os concluintes ( ver Tabela 3.13 ).
Entre os com nota nula, uma pequena minoria ( 1,7 \% ) utilizou a questão para protestar.
Apesar do baixo desempenho, em ambos os grupos a nota máxima de 100,0 pontos foi alcançada.
Os coeficientes de assimetria positivos indicaram uma concentração à direita e espalhamento à esquerda.
O Gráfico 3.11 apresenta a distribuição das notas na questão 39 discursiva do Componente Específico por grupo de estudantes.
A moda se situou no primeiro intervalo, aquele com valores entre 0 e 10 incluindo os extremos do intervalo.
A impressão de assimetria informada pelos coeficientes de assimetria positivos ( Quadro 3.4 ) foi reforçada no gráfico pela concentração no primeiro intervalo.
A distribuição apresentou um outro máximo local no quinto intervalo, de notas razoavelmente baixas, não indicando uma polarização nas habilidades dos alunos. 3.3.2.3 Comentários sobre as respostas à Questão Discursiva 39 Essa questão foi elaborada com o objetivo de avaliar o domínio do conhecimento relativo da Temática Contabilidade Societária, abordando especificamente, conhecimentos acerca do balanço patrimonial e da demonstração do resultado do exercício.
O nível de dificuldades foi médio e atendeu às habilidades e competências descritas.
Com relação ao desempenho dos estudantes e considerando a escala de pontos adotada na correção, variando de 0 ( zero ), 1,75 ; 3,5 ; 5,25 ; e 7,0 verificou -se que entre as provas com respostas válidas ( {\EmojiFont ─} 30 \% do total de provas ), 70 \% dos estudantes acertaram metade ou menos da questão.
Os 30 \% restantes, podem ser distribuídos em {\EmojiFont ─}10 \% com o total de acerto em torno de 20 \% com apenas 1 ; 4 ou menos de acerto.
O domínio do conteúdo pode ser verificado a partir de dois cenários : Cenário 1 – provas com respostas que demonstravam claramente o conhecimento do assunto, dos termos técnicos, do sistema de classificação contábil, entre outros.
Cenário 2 – provas com respostas que apresentavam apenas valores, demonstrando o domínio do conteúdo pelo uso correto dos valores das respectivas rubricas contábeis.
Verificou -se que os erros mais comuns se concentravam em erros de soma dos valores, reclamação pela falta da calculadora e algumas vezes, do excesso de rubricas, de valores, e até demonstração de insatisfação devido à obrigatoriedade da prova.
Algumas respostas foram direcionadas pela incoerência do envolvimento de estudantes ingressantes e concluintes, dando a entender se tratar de universos distintos e de momentos distintos no curso, para se submeterem a uma mesma prova.
Essa questão foi elaborada atendendo às exigências previamente ajustadas e devido à importância do conteúdo abordado na mesma, foi um excelente indicador de avaliação do conhecimento contábil, na medida em que os dois universos de estudantes ( ingressantes e concluintes ) dominavam o assunto abordado na questão. 3.3.2.4 Análise da Questão Discursiva 40 do Componente Específico A questão 40, cujas habilidades aferidas encontram -se descritas no Quadro 3.5, que ainda contém outras informações, seguido da Tabela 3.14 e do Gráfico 3.12, foi aquela em que os grupos de estudantes observados obtiveram pior desempenho entre as questões discursivas do Componente Específico.
A nota média do grupo dos ingressantes foi de 6,64, e a dos concluintes 7,61.
Nesta questão ocorre também o menor percentual de notas nulas : 79,8 \% entre os ingressantes e 78,9 \% entre os concluintes.
Aqui também a maioria das provas com nota zero decorreu de provas deixadas em branco : 76,7 \% entre os ingressantes e 78,5 \% entre os concluintes ( ver Tabela 3.14 ).
Uma pequena minoria ( 1,6 \% ), dentre aqueles com nota nula, utilizou a questão para protestar.
Em ambos os grupos a nota máxima de 100,0 pontos foi alcançada.
Os coeficientes de assimetria positivos indicaram uma inclinação para a esquerda.
O Gráfico 3.12 apresenta a distribuição das notas na questão 40 discursiva do Componente Específico por grupo de estudantes.
A moda se situou no primeiro intervalo, aquele com valores entre 0 e 10 incluindo os extremos do intervalo.
A impressão de assimetria informada pelos coeficientes de assimetria positivos ( Quadro 3.5 ) foi reforçada no gráfico pela concentração no primeiro intervalo.
A distribuição apresentou um outro máximo local no terceiro e quinto intervalo para os ingressantes e concluintes, respectivamente, indicando uma ligeira simetria entre as notas não nulas, mas não polarização das habilidades dos alunos. 3.3.2.5 Comentários sobre as respostas à Questão Discursiva 40 Essa questão foi elaborada com o objetivo de avaliar o domínio do conhecimento relativo à Temática Contabilidade e Análise de Custos, abordando especificamente, conhecimentos acerca dos Métodos de Custeio com aplicação prática em um “ Estudo de Caso Simulado ”, com nível de dificuldade mediano.
Com relação ao desempenho dos estudantes e considerando a escala de pontos adotada na correção, variando de 0 ( zero ), 1,75 ; 3,5 ; 5,25 ; e 7,0, verificou -se que entre as provas com respostas válidas ( – 30 \% do total de provas ), em torno de 80 \% dos estudantes se concentraram em torno da nota 1,75, demonstrando pouco domínio desse conteúdo.
Na mesma tendência, dos 20 \% restantes, menos de 5 \% atingiram a totalidade da questão, com grande concentração entre 3,5 e 5,25.
Observou -se com relação ao domínio do conteúdo, pequena concentração de respostas que demonstravam domínio do conteúdo e uso adequado dos termos técnicos, SMJ, em torno de 5 \%, no máximo, enquanto a grande maioria se limitava a apresentar números e valores, com alguma coerência de leitura e interpretação de conteúdo.
Verificaram -se muitos erros de cálculos e muita reclamação pela falta de calculadora.
No geral, entre as provas com respostas em torno de 47.000, não foi observado fuga ao tema, a não ser para os casos de críticas, ofensas e agressões ao processo e ao sistema ENADE.
É opinião da banca de corretores que essa questão foi elaborada, atendendo às exigências previamente ajustadas, com dificuldade mediana, elaborada na forma de “ Estudo de Caso ”, envolvendo assunto e conteúdo presentes nas Diretrizes Curriculares do Curso, “ Contabilidade e Análise de Custos ”, portanto, um assunto presente, atual e obrigatório para todos os Cursos de Ciências Contábeis do País, mas que demonstrou não ser adequadamente tratado, ou melhor, desenvolvido nos programas do Curso.
O resultado apurado poderia ser justificado pelo fato de se tratar de conteúdo de final de curso e, assim, ser um conteúdo com que os estudantes ingressantes ainda não tiveram contato.
Mas, se observarmos que, nessa área, os estudantes concluintes somam mais do dobro de ingressantes, esse argumento perde valor.

% matched lemmas: bruto, originalidade, valorativo
\end{textsample}
